\section{Synthèse des quatre blocs}

Les projets précédents sont complémentaires : aucun ne cherche à couvrir artificiellement l'ensemble du référentiel. Leur combinaison permet en revanche de montrer une progression allant de l'analyse d'un existant complexe à la conception, au développement, à l'industrialisation puis au pilotage d'un projet collectif.

\subsection*{Bloc 1 --- Planifier et organiser un projet de développement logiciel}

Le bloc~1 repose principalement sur FCU / Ediacara et Lynx Immo. Ediacara démontre l'analyse d'un environnement professionnel complexe, l'identification des risques et la conception d'une adaptation compatible avec la production existante. Lynx Immo complète cette dimension par une étude de faisabilité plus complète, intégrant les aspects techniques, humains et financiers, puis par une feuille de route, des jalons, des outils de suivi et des arbitrages de périmètre. Le homelab apporte un complément sur la veille et l'expérimentation technologique, qui reste à formaliser davantage avec des sources anglophones.

\subsection*{Bloc 2 --- Concevoir et développer des solutions logicielles}

Les preuves sont volontairement réparties entre plusieurs projets. S-123 illustre principalement la conception d'architecture et l'évolution d'un modèle devenu trop couplé à la norme. \texttt{piv2glz} apporte une preuve claire d'intégrité du code par les tests automatisés et la CI. Brise-glace constitue la preuve principale du front-end avec React, TypeScript et une interface temps réel, tandis que FeastVerse démontre le développement back-end avec Spring Boot, PostgreSQL, JWT et une architecture en couches. Le projet Data / IA doit enfin porter C2.5 ; tant que ce chapitre n'est pas finalisé, cette compétence reste le principal point à sécuriser dans ce bloc.

\subsection*{Bloc 3 --- Piloter la mise en production et l'évolution}

Ce bloc s'appuie sur plusieurs niveaux d'industrialisation. \texttt{piv2glz} montre une chaîne simple et reproductible où les tests précèdent la construction des livrables. Ediacara apporte l'expérience d'une mise en production contrôlée dans un environnement legacy, avec validation sur RE7, surveillance des jobs et documentation de la procédure. FeastVerse et Brise-glace illustrent des déploiements automatisés d'applications, tandis que le homelab permet d'expérimenter GitHub Actions, Helm, Argo CD, Kubernetes et une stack d'observabilité. Ces expériences montrent surtout qu'une mise en production fiable combine versionnement, validation, automatisation, surveillance et documentation selon les contraintes du contexte.

\subsection*{Bloc 4 --- Piloter l'équipe du projet de développement logiciel}

Lynx Immo constitue la preuve principale du bloc~4. Mon rôle de chef de projet et de référent technique m'a amené à identifier les compétences nécessaires, clarifier les responsabilités, répartir les tâches, coordonner l'activité, accompagner les membres et réajuster le projet face aux retards et aux variations de disponibilité. Le contexte reste celui d'une équipe étudiante : je n'ai ni recruté l'équipe ni exercé de fonction RH. La montée en compétences s'est faite principalement par documentation, prototypes et échanges internes, et l'évaluation du fonctionnement collectif par rétrospective et entretiens de fin de projet. Ediacara complète ce bloc par la transmission et la capitalisation de connaissances techniques.

\section{Tableau final de couverture}

Le tableau~\ref{tab:synthese-competences} fournit une dernière carte de lecture. Il indique le projet utilisé comme preuve principale et renvoie vers le chapitre où cette preuve est développée.

\begin{longtable}{p{1.4cm} p{5.6cm} p{4.0cm} p{3.4cm}}
\caption{Synthèse des compétences et des projets de référence}\label{tab:synthese-competences}\\
\toprule
\textbf{Comp.} & \textbf{Objet} & \textbf{Projet principal} & \textbf{Renvoi} \\
\midrule
\endfirsthead
\toprule
\textbf{Comp.} & \textbf{Objet} & \textbf{Projet principal} & \textbf{Renvoi} \\
\midrule
\endhead
C1.1 & Étude de faisabilité & Lynx Immo & p.~\pageref{chap:lynx} \\
C1.2 & Veille technologique & Homelab & p.~\pageref{chap:homelab} \\
C1.3 & Conception de la solution & FCU / Ediacara & p.~\pageref{chap:fcu} \\
C1.4 & Feuille de route & Lynx Immo & p.~\pageref{chap:lynx} \\
C1.5 & Mise en place du projet & Lynx Immo & p.~\pageref{chap:lynx} \\
C1.6 & Mesure de la performance & Lynx Immo & p.~\pageref{chap:lynx} \\
C1.7 & Supervision du projet & Lynx Immo & p.~\pageref{chap:lynx} \\
\midrule
C2.1 & Architecture logicielle & S-123 & p.~\pageref{chap:s123} \\
C2.2 & Intégrité du code & \texttt{piv2glz} & p.~\pageref{chap:piv2glz} \\
C2.3 & Développement front-end & Brise-glace & p.~\pageref{chap:brise-glace} \\
C2.4 & Développement back-end & FeastVerse & p.~\pageref{chap:feastverse} \\
C2.5 & Traitement de données & Projet Data / IA & p.~\pageref{chap:data-ia} \\
\midrule
C3.1 & Intégration continue & \texttt{piv2glz} & p.~\pageref{chap:piv2glz} \\
C3.2 & Tests itératifs & \texttt{piv2glz} & p.~\pageref{chap:piv2glz} \\
C3.3 & Surveillance des automatisations & FCU / Ediacara & p.~\pageref{chap:fcu} \\
C3.4 & Déploiement continu & Homelab & p.~\pageref{chap:homelab} \\
C3.5 & Optimisation et observabilité & Homelab & p.~\pageref{chap:homelab} \\
C3.6 & Documentation technique & FCU / Ediacara & p.~\pageref{chap:fcu} \\
\midrule
C4.1 & Compétences nécessaires & Lynx Immo & p.~\pageref{chap:lynx} \\
C4.2 & Constitution de l'équipe & Lynx Immo & p.~\pageref{chap:lynx} \\
C4.3 & Coordination & Lynx Immo & p.~\pageref{chap:lynx} \\
C4.4 & Accompagnement & Lynx Immo & p.~\pageref{chap:lynx} \\
C4.5 & Montée en compétences & Lynx Immo & p.~\pageref{chap:lynx} \\
C4.6 & Évaluation de la performance collective & Lynx Immo & p.~\pageref{chap:lynx} \\
\bottomrule
\end{longtable}

Deux points restent à consolider avant la version finale : C2.5 par la rédaction du chapitre Data / IA, et C1.2 par une veille technologique explicitement sourcée, notamment à partir de références anglophones. La prise en compte du numérique responsable doit également rester visible dans les choix d'architecture et d'exploitation, sans être isolée artificiellement du reste des projets.
